%============================================================
%  Bd. 05 - Äquivalenzrelationen
%============================================================

\documentclass[main.tex]{subfiles}

\ifSubfilesClassLoaded{
  \BandSetupStandalone{B05}
}{
}

\title{Bd. 05 - Äquivalenzrelationen}
\author{Martin Kunze}
\date{}
\setcounter{file}{5}

\begin{document}

\title{Bd. 05 - Äquivalenzrelationen}
\author{Martin Kunze}
\date{}
\hypersetup{pageanchor=false}
\maketitle
\hypersetup{pageanchor=true}
\setcounter{tocdepth}{3}
\tableofcontents
\setcounter{file}{5}
\renewcommand{\FormulaBandID}{5}

\chapter{Einleitung}

Äquivalenzrelationen strukturieren einen Träger durch Klassen
ununterscheidbarer Elemente. Nach der allgemeinen Theorie folgen
Quotientenmengen, die zugehörigen Projektionsabbildungen und
Gleichmächtigkeit als wichtiges funktionsabhängiges Beispiel. Die
Ununterscheidbarkeitsrelation einer Mengenfamilie dient als durchgehendes
konkretes Beispiel.

\chapter{Äquivalenzrelationen und Quotienten}

\section{Begriff der Äquivalenzrelation}

\subsection{Axiome einer Äquivalenzrelation}

\FormulaAxiomDeltaK[Reflexivität]{%
  x\in A \vdash (x,x)\in R%
}{EqRelReflexivity}{
  \DeltaRow{Mengen}{A \dsep R \dsep x}
}

\FormulaAxiomDeltaK[Symmetrie]{%
  (x,y)\in R \vdash (y,x)\in R%
}{EqRelSymmetry}{
  \DeltaRow{Mengen}{A \dsep R \dsep x \dsep y}
}

\FormulaAxiomDeltaK[Transitivität]{%
  (x,y)\in R \dsep (y,z)\in R \vdash (x,z)\in R%
}{EqRelTransitivity}{
  \DeltaRow{Mengen}{A \dsep R \dsep x \dsep y \dsep z}
}

\subsection{Definition der Äquivalenzrelation}

\FormulaDefDeltaK[Begriff der Äquivalenzrelation]{\EqRel{R,A}}{EqRelDef}{
  \DeltaRow{Mengen}{A \dsep R \dsep x \dsep y \dsep z}
  \DeltaRow{\textbf{Axiome}}{}
  \DeltaRow{Relation auf \(A\)}
           {R \subseteq A \times A}
           [\FormulaRefAuto{F \subseteq A \times B}]
  \DeltaRow{Reflexivität}
           {x\in A \vdash (x,x)\in R}
           [\FormulaRefAuto{x\in A \vdash (x,x)\in R}]
  \DeltaRow{Symmetrie}
           {(x,y)\in R \vdash (y,x)\in R}
           [\FormulaRefAuto{(x,y)\in R \vdash (y,x)\in R}]
  \DeltaRow{Transitivität}
           {(x,y)\in R \dsep (y,z)\in R \vdash (x,z)\in R}
           [\FormulaRefAuto{(x,y)\in R \dsep (y,z)\in R \vdash (x,z)\in R}]
  \DeltaRow{\textbf{Neue Symbole}}{}
  \DeltaPrem{Äquivalenzrelationen}{}
}

\subsection{Beispiele von Äquivalenzrelationen}

\subsubsection{Ununterscheidbarkeitsrelation einer Mengenfamilie}

Ein wichtiges Beispiel entsteht aus einer Mengenfamilie \(M\). In der
Terminologie der Rough-Set-Theorie ist dies eine
Ununterscheidbarkeitsrelation (indiscernibility relation): Zwei Elemente des
Trägers \(\bigcup M\) sind bezüglich \(M\) ununterscheidbar, wenn sie in jedem
Mitglied von \(M\) entweder beide enthalten oder beide nicht enthalten sind.

\FormulaDefDeltaK[Ununterscheidbarkeitsrelation einer Mengenfamilie]{%
  \SameOccRel{M}\coloneqq
  \{(a,b)\in \bigcup M\times\bigcup M\mid
  \forall X\in M\,(a\in X\leftrightarrow b\in X)\}
}{SameOccRelDef}{%
  \DeltaRow{Mengen}{M\dsep X\dsep a\dsep b}
  \DeltaRow{Relationen}{\SameOccRel{M}}
  \DeltaRow{Neue Symbole}{\SameOccRel{M}}
}

\FormulaThmDeltaK[Charakterisierung der Ununterscheidbarkeitsrelation]{%
  (a,b)\in \SameOccRel{M}
  \eqvdash
  a\in\bigcup M\land b\in\bigcup M\land
  \forall X\in M\,(a\in X\leftrightarrow b\in X)
}{SameOccRelChar}{%
  \DeltaRow{Mengen}{M\dsep X\dsep a\dsep b}
  \DeltaRow{Relationen}{\SameOccRel{M}}
}
\begin{tabproofsplit}
\proofpart{\(\vdash\)}
  \proofstep{1}{(a,b)\in \SameOccRel{M}}{\rA}
  \proofstep{}{%
    \SameOccRel{M}\coloneqq
    \{(a,b)\in \bigcup M\times\bigcup M\mid
    \forall X\in M\,(a\in X\leftrightarrow b\in X)\}}{%
    \FormulaRefAuto{SameOccRelDef}}
  \proofstep{1}{(a,b)\in\bigcup M\times\bigcup M\land
    \forall X\in M\,(a\in X\leftrightarrow b\in X)}{%
    \FormulaRefAuto{x \in \{u \in A \mid P(u)\} \eqvdash x \in A \land P(x)}{2,1}}
  \proofstep{1}{(a,b)\in\bigcup M\times\bigcup M}{\rAEa{3}}
  \proofstep{1}{\forall X\in M\,(a\in X\leftrightarrow b\in X)}{\rAEb{3}}
  \proofstep{1}{a\in\bigcup M}{\FormulaRefAuto{(a,b)\in A\times B\vdash a\in A}{4}}
  \proofstep{1}{b\in\bigcup M}{\FormulaRefAuto{(a,b)\in A\times B\vdash b\in B}{4}}
  \proofstep{1}{a\in\bigcup M\land b\in\bigcup M\land
    \forall X\in M\,(a\in X\leftrightarrow b\in X)}{\rAI{6,\rAI{7,5}}}
\closeproofpart

\proofpart{\(\dashv\)}
  \proofstep{1}{a\in\bigcup M\land b\in\bigcup M\land
    \forall X\in M\,(a\in X\leftrightarrow b\in X)}{\rA}
  \proofstep{}{%
    \SameOccRel{M}\coloneqq
    \{(a,b)\in \bigcup M\times\bigcup M\mid
    \forall X\in M\,(a\in X\leftrightarrow b\in X)\}}{%
    \FormulaRefAuto{SameOccRelDef}}
  \proofstep{1}{a\in\bigcup M}{\rAEa{1}}
  \proofstep{1}{b\in\bigcup M}{\rAEa{\rAEb{1}}}
  \proofstep{1}{\forall X\in M\,(a\in X\leftrightarrow b\in X)}{\rAEb{\rAEb{1}}}
  \proofstep{1}{(a,b)\in\bigcup M\times\bigcup M}{%
    \FormulaRefAuto{a\in A\dsep b\in B\vdash (a,b)\in A\times B}{3,4}}
  \proofstep{1}{(a,b)\in\bigcup M\times\bigcup M\land
    \forall X\in M\,(a\in X\leftrightarrow b\in X)}{\rAI{6,5}}
  \proofstep{1}{(a,b)\in \SameOccRel{M}}{%
    \FormulaRefAuto{x \in \{u \in A \mid P(u)\} \eqvdash x \in A \land P(x)}{2,7}}
\closeproofpart
\end{tabproofsplit}

\FormulaThmDeltaK[Trägerbedingung der Ununterscheidbarkeitsrelation]{%
  \SameOccRel{M}\subseteq\bigcup M\times\bigcup M
}{SameOccRelCarrier}{%
  \DeltaRow{Mengen}{M\dsep X\dsep a\dsep b\dsep c}
  \DeltaRow{Relationen}{\SameOccRel{M}}
}
\begin{tabproofwide}
  \proofstepwidestar[1]{(a,b)\in \SameOccRel{M}}{\rA}
  \proofstepwidestar[1]{a\in\bigcup M\land b\in\bigcup M\land
    \forall X\in M\,(a\in X\leftrightarrow b\in X)}{%
    \FormulaRefAuto{SameOccRelChar}{1}}
  \proofstepwidestar[1]{a\in\bigcup M}{\rAEa{2}}
  \proofstepwidestar[1]{b\in\bigcup M}{\rAEa{\rAEb{2}}}
  \proofstepwidestar[1]{(a,b)\in\bigcup M\times\bigcup M}{%
    \FormulaRefAuto{a\in A\dsep b\in B\vdash (a,b)\in A\times B}{3,4}}
  \proofstepwidestar[]{\SameOccRel{M}\subseteq\bigcup M\times\bigcup M}{%
    \FormulaRefAuto{A\subseteq B \coloneq
    \forall x\,(x\in A \rightarrow x\in B)}{\rUI{\rRI{1,5}}}}
\end{tabproofwide}

\FormulaThmDeltaK[Reflexivität der Ununterscheidbarkeitsrelation]{%
  a\in\bigcup M \vdash (a,a)\in \SameOccRel{M}
}{SameOccRelReflexivity}{%
  \DeltaRow{Mengen}{M\dsep X\dsep a}
  \DeltaRow{Relationen}{\SameOccRel{M}}
}
\begin{tabproof}
  \proofstep{1}{a\in\bigcup M}{\rA}
  \proofstep{2}{X\in M}{\rA}
  \proofstep{1,2}{a\in X\leftrightarrow a\in X}{\FormulaRefAuto{P\leftrightarrow P}}
  \proofstep{1}{\forall X\in M\,(a\in X\leftrightarrow a\in X)}{\rUI{\rRI{2,3}}}
  \proofstep{1}{a\in\bigcup M\land a\in\bigcup M\land
    \forall X\in M\,(a\in X\leftrightarrow a\in X)}{\rAI{1,\rAI{1,4}}}
  \proofstep{1}{(a,a)\in \SameOccRel{M}}{%
    \FormulaRefAuto{SameOccRelChar}{5}}
\end{tabproof}

\FormulaThmDeltaK[Symmetrie der Ununterscheidbarkeitsrelation]{%
  (a,b)\in \SameOccRel{M} \vdash (b,a)\in \SameOccRel{M}
}{SameOccRelSymmetry}{%
  \DeltaRow{Mengen}{M\dsep X\dsep a\dsep b}
  \DeltaRow{Relationen}{\SameOccRel{M}}
}
\begin{tabproofwide}
  \proofstepwidestar[1]{(a,b)\in \SameOccRel{M}}{\rA}
  \proofstepwidestar[1]{a\in\bigcup M\land b\in\bigcup M\land
    \forall X\in M\,(a\in X\leftrightarrow b\in X)}{%
    \FormulaRefAuto{SameOccRelChar}{1}}
  \proofstepwidestar[1]{b\in\bigcup M}{\rAEa{\rAEb{2}}}
  \proofstepwidestar[1]{a\in\bigcup M}{\rAEa{2}}
  \proofstepwidestar[1]{\forall X\in M\,(b\in X\leftrightarrow a\in X)}{%
    \FormulaRefAuto{\forall x(P(x)\leftrightarrow Q(x))\vdash
    \forall x(Q(x)\leftrightarrow P(x))}{\rAEb{\rAEb{2}}}}
  \proofstepwidestar[1]{b\in\bigcup M\land a\in\bigcup M\land
    \forall X\in M\,(b\in X\leftrightarrow a\in X)}{\rAI{3,\rAI{4,5}}}
  \proofstepwidestar[1]{(b,a)\in \SameOccRel{M}}{%
    \FormulaRefAuto{SameOccRelChar}{6}}
\end{tabproofwide}

\FormulaThmDeltaK[Transitivität der Ununterscheidbarkeitsrelation]{%
  (a,b)\in \SameOccRel{M}\dsep (b,c)\in \SameOccRel{M}
  \vdash (a,c)\in \SameOccRel{M}
}{SameOccRelTransitivity}{%
  \DeltaRow{Mengen}{M\dsep X\dsep a\dsep b\dsep c}
  \DeltaRow{Relationen}{\SameOccRel{M}}
}
\begin{tabproofwide}
  \proofstepwidestar[1]{(a,b)\in \SameOccRel{M}}{\rA}
  \proofstepwidestar[2]{(b,c)\in \SameOccRel{M}}{\rA}
  \proofstepwidestar[1]{a\in\bigcup M\land b\in\bigcup M\land
    \forall X\in M\,(a\in X\leftrightarrow b\in X)}{%
    \FormulaRefAuto{SameOccRelChar}{1}}
  \proofstepwidestar[2]{b\in\bigcup M\land c\in\bigcup M\land
    \forall X\in M\,(b\in X\leftrightarrow c\in X)}{%
    \FormulaRefAuto{SameOccRelChar}{2}}
  \proofstepwidestar[1]{a\in\bigcup M}{\rAEa{3}}
  \proofstepwidestar[2]{c\in\bigcup M}{\rAEa{\rAEb{4}}}
  \proofstepwidestar[1,2]{\forall X\in M\,(a\in X\leftrightarrow c\in X)}{%
    \FormulaRefAuto{\forall x(P(x)\leftrightarrow Q(x)),
    \forall x(Q(x)\leftrightarrow R(x))\vdash
    \forall x(P(x)\leftrightarrow R(x))}{\rAEb{\rAEb{3}},\rAEb{\rAEb{4}}}}
  \proofstepwidestar[1,2]{a\in\bigcup M\land c\in\bigcup M\land
    \forall X\in M\,(a\in X\leftrightarrow c\in X)}{\rAI{5,\rAI{6,7}}}
  \proofstepwidestar[1,2]{(a,c)\in \SameOccRel{M}}{%
    \FormulaRefAuto{SameOccRelChar}{8}}
\end{tabproofwide}

\FormulaThmDeltaK[Ununterscheidbarkeitsrelation ist Äquivalenzrelation]{%
  \EqRel{\SameOccRel{M},\bigcup M}
}{SameOccRelEqRel}{%
  \DeltaRow{Mengen}{M\dsep X\dsep a\dsep b\dsep c}
  \DeltaRow{Relationen}{\SameOccRel{M}}
}
\begin{tabproof}
  \proofstep{}{\SameOccRel{M}\subseteq\bigcup M\times\bigcup M}{%
    \FormulaRefAuto{SameOccRelCarrier}}
  \proofstep{}{a\in\bigcup M \rightarrow (a,a)\in \SameOccRel{M}}{%
    \FormulaRefAuto{SameOccRelReflexivity}}
  \proofstep{}{(a,b)\in \SameOccRel{M} \rightarrow (b,a)\in \SameOccRel{M}}{%
    \FormulaRefAuto{SameOccRelSymmetry}}
  \proofstep{}{%
    \bigl((a,b)\in \SameOccRel{M}\land (b,c)\in \SameOccRel{M}\bigr)
    \rightarrow (a,c)\in \SameOccRel{M}}{%
    \FormulaRefAuto{SameOccRelTransitivity}}
  \proofstep{}{\EqRel{\SameOccRel{M},\bigcup M}}{%
    \FormulaRefAuto{EqRelDef}{1,2,3,4}}
\end{tabproof}

\subsection{Äquivalenzklassen}

\FormulaDefDeltaK[Äquivalenzklasse eines Elements]{%
  \EqRel{R,A}\dsep x\in A\vdash
  \EqClass{x}{R}{A}\coloneqq
  \{y\in A\mid (x,y)\in R\}
}{EqClassDef}{%
  \DeltaRow{Mengen}{A\dsep R\dsep x\dsep y}
  \DeltaRow{Relationen}{R}
  \DeltaRow{Neue Symbole}{\EqClass{x}{R}{A}}
}

\FormulaThmDeltaK[Charakterisierung der Äquivalenzklasse]{%
  \EqRel{R,A}\dsep x\in A
  \vdash
  y\in \EqClass{x}{R}{A}
  \leftrightarrow
  y\in A\land (x,y)\in R
}{EqClassChar}{%
  \DeltaRow{Mengen}{A\dsep R\dsep x\dsep y}
  \DeltaRow{Relationen}{R}
}
\begin{tabproof}
  \proofstep{1}{\EqRel{R,A}}{\rA}
  \proofstep{2}{x\in A}{\rA}
  \proofstep{1,2}{%
    \EqClass{x}{R}{A}\coloneqq
    \{y\in A\mid (x,y)\in R\}}{%
    \FormulaRefAuto{EqClassDef}{1,2}}

  \proofstep{4}{y\in \EqClass{x}{R}{A}}{\rA}
  \proofstep{1,2,4}{y\in A\land (x,y)\in R}{%
    \FormulaRefAuto{x \in \{u \in A \mid P(u)\} \eqvdash x \in A \land P(x)}{3,4}}
  \proofstep{1,2}{%
    y\in \EqClass{x}{R}{A}
    \rightarrow y\in A\land (x,y)\in R}{\rRI{4,5}}

  \proofstep{7}{y\in A\land (x,y)\in R}{\rA}
  \proofstep{1,2,7}{y\in \EqClass{x}{R}{A}}{%
    \FormulaRefAuto{x \in \{u \in A \mid P(u)\} \eqvdash x \in A \land P(x)}{3,7}}
  \proofstep{1,2}{%
    y\in A\land (x,y)\in R
    \rightarrow y\in \EqClass{x}{R}{A}}{\rRI{7,8}}

  \proofstep{1,2}{%
    y\in \EqClass{x}{R}{A}
    \leftrightarrow
    y\in A\land (x,y)\in R}{\rLRI{6,9}}
\end{tabproof}

\FormulaThmDeltaK[Repräsentant liegt in seiner Äquivalenzklasse]{%
  \EqRel{R,A}\dsep x\in A
  \vdash
  x\in \EqClass{x}{R}{A}
}{EqClassRepresentative}{%
  \DeltaRow{Mengen}{A\dsep R\dsep x}
  \DeltaRow{Relationen}{R}
}
\begin{tabproof}
  \proofstep{1}{\EqRel{R,A}}{\rA}
  \proofstep{2}{x\in A}{\rA}
  \proofstep{1,2}{(x,x)\in R}{\FormulaRefAuto{EqRelDef}{1,2}}
  \proofstep{1,2}{x\in A\land (x,x)\in R}{\rAI{2,3}}
  \proofstep{1,2}{x\in \EqClass{x}{R}{A}\leftrightarrow
    x\in A\land (x,x)\in R}{\FormulaRefAuto{EqClassChar}{1,2}}
  \proofstep{1,2}{x\in \EqClass{x}{R}{A}}{%
    \FormulaRefAuto{P \leftrightarrow Q, Q \vdash P}{5,4}}
\end{tabproof}

\FormulaThmDeltaK[Elemente einer Äquivalenzklasse liegen im Grundbereich]{%
  \EqRel{R,A}\dsep x\in A\dsep y\in \EqClass{x}{R}{A}
  \vdash y\in A
}{EqClassElementInCarrier}{%
  \DeltaRow{Mengen}{A\dsep R\dsep x\dsep y}
  \DeltaRow{Relationen}{R}
}
\begin{tabproof}
  \proofstep{1}{\EqRel{R,A}}{\rA}
  \proofstep{2}{x\in A}{\rA}
  \proofstep{3}{y\in \EqClass{x}{R}{A}}{\rA}
  \proofstep{1,2}{y\in \EqClass{x}{R}{A}\leftrightarrow
    y\in A\land (x,y)\in R}{\FormulaRefAuto{EqClassChar}{1,2}}
  \proofstep{1,2,3}{y\in A\land (x,y)\in R}{%
    \FormulaRefAuto{P \leftrightarrow Q, P \vdash Q}{4,3}}
  \proofstep{1,2,3}{y\in A}{\rAEa{5}}
\end{tabproof}

\FormulaThmDeltaK[Elemente einer Äquivalenzklasse sind äquivalent zum Repräsentanten]{%
  \EqRel{R,A}\dsep x\in A\dsep y\in \EqClass{x}{R}{A}
  \vdash (x,y)\in R
}{EqClassElementRelated}{%
  \DeltaRow{Mengen}{A\dsep R\dsep x\dsep y}
  \DeltaRow{Relationen}{R}
}
\begin{tabproof}
  \proofstep{1}{\EqRel{R,A}}{\rA}
  \proofstep{2}{x\in A}{\rA}
  \proofstep{3}{y\in \EqClass{x}{R}{A}}{\rA}
  \proofstep{1,2}{y\in \EqClass{x}{R}{A}\leftrightarrow
    y\in A\land (x,y)\in R}{\FormulaRefAuto{EqClassChar}{1,2}}
  \proofstep{1,2,3}{y\in A\land (x,y)\in R}{%
    \FormulaRefAuto{P \leftrightarrow Q, P \vdash Q}{4,3}}
  \proofstep{1,2,3}{(x,y)\in R}{\rAEb{5}}
\end{tabproof}

\FormulaThmDeltaK[Äquivalente Elemente haben dieselbe Äquivalenzklasse]{%
  \EqRel{R,A}\dsep x\in A\dsep y\in A\dsep (x,y)\in R
  \vdash
  \EqClass{x}{R}{A}=\EqClass{y}{R}{A}
}{EqClassEqualFromRel}{%
  \DeltaRow{Mengen}{A\dsep R\dsep x\dsep y\dsep z}
  \DeltaRow{Relationen}{R}
}
\begin{tabproofsplitwide}
  \proofpartwideindR[Teilmengenrichtung \(\EqClass{x}{R}{A}\subseteq\EqClass{y}{R}{A}\)]{%
    \EqRel{R,A}\dsep x\in A\dsep y\in A\dsep (x,y)\in R
    \vdash
    \EqClass{x}{R}{A}\subseteq\EqClass{y}{R}{A}
  }{%
    \EqRel{R,A}\dsep x\in A\dsep y\in A\dsep (x,y)\in R
    \vdash
    \EqClass{x}{R}{A}\subseteq\EqClass{y}{R}{A}
  }[EqClassEqualFromRelSubsetLeft]
    \proofstepwidestar[1]{\EqRel{R,A}}{\rA}
    \proofstepwidestar[2]{x\in A}{\rA}
    \proofstepwidestar[3]{y\in A}{\rA}
    \proofstepwidestar[4]{(x,y)\in R}{\rA}
    \proofstepwidestar[5]{z\in \EqClass{x}{R}{A}}{\rA}
    \proofstepwidestar[1,2,5]{z\in A}{%
      \FormulaRefAuto{EqClassElementInCarrier}{1,2,5}}
    \proofstepwidestar[1,2,5]{(x,z)\in R}{%
      \FormulaRefAuto{EqClassElementRelated}{1,2,5}}
    \proofstepwidestar[1,4]{(y,x)\in R}{\FormulaRefAuto{EqRelSymmetry}{4}}
    \proofstepwidestar[1,2,4,5]{(y,z)\in R}{%
      \FormulaRefAuto{(x,y)\in R \dsep (y,z)\in R \vdash (x,z)\in R}[ax]{8,7}}
    \proofstepwidestar[1,2,3,4,5]{z\in A\land (y,z)\in R}{\rAI{6,9}}
    \proofstepwidestar[1,3]{z\in \EqClass{y}{R}{A}\leftrightarrow
      z\in A\land (y,z)\in R}{\FormulaRefAuto{EqClassChar}{1,3}}
    \proofstepwidestar[1,2,3,4,5]{z\in \EqClass{y}{R}{A}}{%
      \FormulaRefAuto{P \leftrightarrow Q, Q \vdash P}{11,10}}
    \proofstepwidestar[1,2,3,4]{\EqClass{x}{R}{A}\subseteq\EqClass{y}{R}{A}}{%
      \FormulaRefAuto{A\subseteq B \coloneq
      \forall x\,(x\in A \rightarrow x\in B)}{\rUI{\rRI{5,12}}}}
  \closeproofpartwideind

  \proofpartwideindR[Teilmengenrichtung \(\EqClass{y}{R}{A}\subseteq\EqClass{x}{R}{A}\)]{%
    \EqRel{R,A}\dsep x\in A\dsep y\in A\dsep (x,y)\in R
    \vdash
    \EqClass{y}{R}{A}\subseteq\EqClass{x}{R}{A}
  }{%
    \EqRel{R,A}\dsep x\in A\dsep y\in A\dsep (x,y)\in R
    \vdash
    \EqClass{y}{R}{A}\subseteq\EqClass{x}{R}{A}
  }[EqClassEqualFromRelSubsetRight]
    \proofstepwidestar[1]{\EqRel{R,A}}{\rA}
    \proofstepwidestar[2]{x\in A}{\rA}
    \proofstepwidestar[3]{y\in A}{\rA}
    \proofstepwidestar[4]{(x,y)\in R}{\rA}
    \proofstepwidestar[5]{z\in \EqClass{y}{R}{A}}{\rA}
    \proofstepwidestar[1,3,5]{z\in A}{%
      \FormulaRefAuto{EqClassElementInCarrier}{1,3,5}}
    \proofstepwidestar[1,3,5]{(y,z)\in R}{%
      \FormulaRefAuto{EqClassElementRelated}{1,3,5}}
    \proofstepwidestar[1,3,4,5]{(x,z)\in R}{%
      \FormulaRefAuto{(x,y)\in R \dsep (y,z)\in R \vdash (x,z)\in R}[ax]{4,7}}
    \proofstepwidestar[1,2,3,4,5]{z\in A\land (x,z)\in R}{\rAI{6,8}}
    \proofstepwidestar[1,2]{z\in \EqClass{x}{R}{A}\leftrightarrow
      z\in A\land (x,z)\in R}{\FormulaRefAuto{EqClassChar}{1,2}}
    \proofstepwidestar[1,2,3,4,5]{z\in \EqClass{x}{R}{A}}{%
      \FormulaRefAuto{P \leftrightarrow Q, Q \vdash P}{10,9}}
    \proofstepwidestar[1,2,3,4]{\EqClass{y}{R}{A}\subseteq\EqClass{x}{R}{A}}{%
      \FormulaRefAuto{A\subseteq B \coloneq
      \forall x\,(x\in A \rightarrow x\in B)}{\rUI{\rRI{5,11}}}}
  \closeproofpartwideind

  \proofpartwide{Schluss}
    \proofstepwidestar[1]{\EqRel{R,A}}{\rA}
    \proofstepwidestar[2]{x\in A}{\rA}
    \proofstepwidestar[3]{y\in A}{\rA}
    \proofstepwidestar[4]{(x,y)\in R}{\rA}
    \proofstepwidestar[1,2,3,4]{\EqClass{x}{R}{A}\subseteq\EqClass{y}{R}{A}}{%
      \FormulaRefAuto{EqClassEqualFromRelSubsetLeft}{1,2,3,4}}
    \proofstepwidestar[1,2,3,4]{\EqClass{y}{R}{A}\subseteq\EqClass{x}{R}{A}}{%
      \FormulaRefAuto{EqClassEqualFromRelSubsetRight}{1,2,3,4}}
    \proofstepwidestar[1,2,3,4]{\EqClass{x}{R}{A}=\EqClass{y}{R}{A}}{%
      \FormulaRefAuto{A \subseteq B, B \subseteq A \vdash A = B}{5,6}}
  \closeproofpartwide
\end{tabproofsplitwide}

\FormulaThmDeltaK[Gleiche Äquivalenzklassen liefern Äquivalenz]{%
  \EqRel{R,A}\dsep x\in A\dsep y\in A\dsep
  \EqClass{x}{R}{A}=\EqClass{y}{R}{A}
  \vdash
  (x,y)\in R
}{EqClassRelFromEqual}{%
  \DeltaRow{Mengen}{A\dsep R\dsep x\dsep y}
  \DeltaRow{Relationen}{R}
}
\begin{tabproof}
  \proofstep{1}{\EqRel{R,A}}{\rA}
  \proofstep{2}{x\in A}{\rA}
  \proofstep{3}{y\in A}{\rA}
  \proofstep{4}{\EqClass{x}{R}{A}=\EqClass{y}{R}{A}}{\rA}
  \proofstep{1,2}{x\in \EqClass{x}{R}{A}}{%
    \FormulaRefAuto{EqClassRepresentative}{1,2}}
  \proofstep{1,2,4}{x\in \EqClass{y}{R}{A}}{\rIE{4,5}}
  \proofstep{1,2,3,4}{(y,x)\in R}{%
    \FormulaRefAuto{EqClassElementRelated}{1,3,6}}
  \proofstep{1,2,3,4}{(x,y)\in R}{\FormulaRefAuto{EqRelSymmetry}{7}}
\end{tabproof}

\FormulaThmDeltaK[Charakterisierung gleicher Äquivalenzklassen]{%
  \EqRel{R,A}\dsep x\in A\dsep y\in A
  \vdash
  \bigl(\EqClass{x}{R}{A}=\EqClass{y}{R}{A}\bigr)
  \leftrightarrow
  (x,y)\in R
}{EqClassEqualIffRel}{%
  \DeltaRow{Mengen}{A\dsep R\dsep x\dsep y}
  \DeltaRow{Relationen}{R}
}
\begin{tabproof}
  \proofstep{1}{\EqRel{R,A}}{\rA}
  \proofstep{2}{x\in A}{\rA}
  \proofstep{3}{y\in A}{\rA}

  \proofstep{4}{\EqClass{x}{R}{A}=\EqClass{y}{R}{A}}{\rA}
  \proofstep{1,2,3,4}{(x,y)\in R}{\FormulaRefAuto{EqClassRelFromEqual}{1,2,3,4}}
  \proofstep{1,2,3}{%
    \bigl(\EqClass{x}{R}{A}=\EqClass{y}{R}{A}\bigr)
    \rightarrow (x,y)\in R}{\rRI{4,5}}

  \proofstep{7}{(x,y)\in R}{\rA}
  \proofstep{1,2,3,7}{\EqClass{x}{R}{A}=\EqClass{y}{R}{A}}{%
    \FormulaRefAuto{EqClassEqualFromRel}{1,2,3,7}}
  \proofstep{1,2,3}{%
    (x,y)\in R
    \rightarrow \EqClass{x}{R}{A}=\EqClass{y}{R}{A}}{\rRI{7,8}}

  \proofstep{1,2,3}{%
    \bigl(\EqClass{x}{R}{A}=\EqClass{y}{R}{A}\bigr)
    \leftrightarrow (x,y)\in R}{\rLRI{6,9}}
\end{tabproof}

\subsection{Quotientenmengen}

\FormulaDefDeltaK[Quotientenmenge einer Äquivalenzrelation]{%
  \EqRel{R,A}\vdash
  \QuotientSet{A}{R}\coloneqq
  \{C\in\powerset(A)\mid
  \exists x\in A\,C=\EqClass{x}{R}{A}\}
}{QuotientSetDef}{%
  \DeltaRow{Mengen}{A\dsep R\dsep C\dsep x}
  \DeltaRow{Relationen}{R}
  \DeltaRow{Äquivalenzklassen}{\EqClass{x}{R}{A}}
  \DeltaRow{Neue Symbole}{\QuotientSet{A}{R}}
}

\FormulaThmDeltaK[Charakterisierung der Quotientenmenge]{%
  \EqRel{R,A}
  \vdash
  C\in \QuotientSet{A}{R}
  \leftrightarrow
  C\in\powerset(A)\land
  \exists x\in A\,C=\EqClass{x}{R}{A}
}{QuotientSetChar}{%
  \DeltaRow{Mengen}{A\dsep R\dsep C\dsep x}
  \DeltaRow{Relationen}{R}
  \DeltaRow{Äquivalenzklassen}{\EqClass{x}{R}{A}}
}
\begin{tabproof}
  \proofstep{1}{\EqRel{R,A}}{\rA}
  \proofstep{1}{%
    \QuotientSet{A}{R}\coloneqq
    \{C\in\powerset(A)\mid
    \exists x\in A\,C=\EqClass{x}{R}{A}\}}{%
    \FormulaRefAuto{QuotientSetDef}{1}}

  \proofstep{3}{C\in \QuotientSet{A}{R}}{\rA}
  \proofstep{1,3}{C\in\powerset(A)\land
    \exists x\in A\,C=\EqClass{x}{R}{A}}{%
    \FormulaRefAuto{x \in \{u \in A \mid P(u)\} \eqvdash x \in A \land P(x)}{2,3}}
  \proofstep{1}{C\in \QuotientSet{A}{R}
    \rightarrow C\in\powerset(A)\land
    \exists x\in A\,C=\EqClass{x}{R}{A}}{\rRI{3,4}}

  \proofstep{6}{C\in\powerset(A)\land
    \exists x\in A\,C=\EqClass{x}{R}{A}}{\rA}
  \proofstep{1,6}{C\in \QuotientSet{A}{R}}{%
    \FormulaRefAuto{x \in \{u \in A \mid P(u)\} \eqvdash x \in A \land P(x)}{2,6}}
  \proofstep{1}{C\in\powerset(A)\land
    \exists x\in A\,C=\EqClass{x}{R}{A}
    \rightarrow C\in \QuotientSet{A}{R}}{\rRI{6,7}}

  \proofstep{1}{C\in \QuotientSet{A}{R}
    \leftrightarrow
    C\in\powerset(A)\land
    \exists x\in A\,C=\EqClass{x}{R}{A}}{\rLRI{5,8}}
\end{tabproof}

\FormulaThmDeltaK[Äquivalenzklasse liegt in der Quotientenmenge]{%
  \EqRel{R,A}\dsep x\in A
  \vdash
  \EqClass{x}{R}{A}\in \QuotientSet{A}{R}
}{EqClassInQuotient}{%
  \DeltaRow{Mengen}{A\dsep R\dsep x\dsep y}
  \DeltaRow{Relationen}{R}
  \DeltaRow{Äquivalenzklassen}{\EqClass{x}{R}{A}}
}
\begin{tabproof}
  \proofstep{1}{\EqRel{R,A}}{\rA}
  \proofstep{2}{x\in A}{\rA}
  \proofstep{3}{y\in \EqClass{x}{R}{A}}{\rA}
  \proofstep{1,2,3}{y\in A}{%
    \FormulaRefAuto{EqClassElementInCarrier}{1,2,3}}
  \proofstep{1,2}{\EqClass{x}{R}{A}\subseteq A}{%
    \FormulaRefAuto{A\subseteq B \coloneq
    \forall x\,(x\in A \rightarrow x\in B)}{\rUI{\rRI{3,4}}}}
  \proofstep{1,2}{\EqClass{x}{R}{A}\in\powerset(A)}{%
    \FormulaRefAuto{x \in \powerset(A)\;\eqvdash\; x \subseteq A}{5}}
  \proofstep{}{ \EqClass{x}{R}{A}=\EqClass{x}{R}{A}}{\rII}
  \proofstep{2}{\exists z\in A\,\EqClass{x}{R}{A}=\EqClass{z}{R}{A}}{\rEI{\rAI{2,7}}}
  \proofstep{1,2}{%
    \EqClass{x}{R}{A}\in\powerset(A)\land
    \exists z\in A\,\EqClass{x}{R}{A}=\EqClass{z}{R}{A}}{\rAI{6,8}}
  \proofstep{1,2}{\EqClass{x}{R}{A}\in \QuotientSet{A}{R}}{%
    \FormulaRefAuto{QuotientSetChar}{1,9}}
\end{tabproof}

\FormulaThmDelta[Elemente der Quotientenmenge sind Teilmengen]{%
  \EqRel{R,A}\dsep C\in \QuotientSet{A}{R}
  \vdash
  C\subseteq A
}{%
  \DeltaRow{Mengen}{A\dsep R\dsep C}
  \DeltaRow{Relationen}{R}
}
\begin{tabproof}
  \proofstep{1}{\EqRel{R,A}}{\rA}
  \proofstep{2}{C\in \QuotientSet{A}{R}}{\rA}
  \proofstep{1,2}{%
    C\in\powerset(A)\land
    \exists x\in A\,C=\EqClass{x}{R}{A}}{%
    \FormulaRefAuto{QuotientSetChar}{1,2}}
  \proofstep{1,2}{C\in\powerset(A)}{\rAEa{3}}
  \proofstep{1,2}{C\subseteq A}{%
    \FormulaRefAuto{x \in \powerset(A)\;\eqvdash\; x \subseteq A}{4}}
\end{tabproof}

\FormulaThmDeltaK[Elemente der Quotientenmenge haben einen Repräsentanten]{%
  \EqRel{R,A}\dsep C\in \QuotientSet{A}{R}
  \vdash
  \exists x\in A\,C=\EqClass{x}{R}{A}
}{QuotientRepresentative}{%
  \DeltaRow{Mengen}{A\dsep R\dsep C\dsep x}
  \DeltaRow{Relationen}{R}
}
\begin{tabproof}
  \proofstep{1}{\EqRel{R,A}}{\rA}
  \proofstep{2}{C\in \QuotientSet{A}{R}}{\rA}
  \proofstep{1,2}{%
    C\in\powerset(A)\land
    \exists x\in A\,C=\EqClass{x}{R}{A}}{%
    \FormulaRefAuto{QuotientSetChar}{1,2}}
  \proofstep{1,2}{\exists x\in A\,C=\EqClass{x}{R}{A}}{\rAEb{3}}
\end{tabproof}


\section{Quotientenprojektionen als Funktionen}

\subsection{Quotientenprojektionen und Quotientenbildabbildungen}

\FormulaDefDeltaK[Funktionsvorschrift der Quotientenprojektion]{%
  \EqRel{R,A}\dsep x\in A\vdash
  \tQuotProj{A}{R}(x)\coloneqq \EqClass{x}{R}{A}
}{QuotProjRuleDef}{%
  \DeltaRow{Mengen}{A\dsep R\dsep x}
  \DeltaRow{Relationen}{R}
  \DeltaRow{Äquivalenzklassen}{\EqClass{x}{R}{A}}
  \DeltaRow{Neue Symbole}{\tQuotProj{A}{R}}
}

\FormulaThmDeltaK[Typisierung der Quotientenprojektionsvorschrift]{%
  \EqRel{R,A}\dsep x\in A\vdash
  \tQuotProj{A}{R}(x)\in\QuotientSet{A}{R}
}{QuotProjRuleInQuotient}{%
  \DeltaRow{Mengen}{A\dsep R\dsep x}
  \DeltaRow{Relationen}{R}
  \DeltaRow{Quotientenmenge}{\QuotientSet{A}{R}}
}
\begin{tabproof}
  \proofstep{1}{\EqRel{R,A}}{\rA}
  \proofstep{2}{x\in A}{\rA}
  \proofstep{1,2}{\tQuotProj{A}{R}(x)=\EqClass{x}{R}{A}}{%
    \FormulaRefAuto{QuotProjRuleDef}{1,2}}
  \proofstep{1,2}{\EqClass{x}{R}{A}\in\QuotientSet{A}{R}}{%
    \FormulaRefAuto{EqClassInQuotient}{1,2}}
  \proofstep{1,2}{\tQuotProj{A}{R}(x)\in\QuotientSet{A}{R}}{\rIE{3,4}}
\end{tabproof}

\FormulaThmDeltaK[Existenz der Quotientenprojektion]{%
  \EqRel{R,A}\vdash
  \exists! q\colon A\to\QuotientSet{A}{R}\,
  \forall x\in A\,q(x)=\tQuotProj{A}{R}(x)
}{QuotProjExists}{%
  \DeltaRow{Mengen}{A\dsep R\dsep x}
  \DeltaRow{Funktionen}{q}
  \DeltaRow{Funktionensymbol}{\tQuotProj{A}{R}}
}
\begin{tabproof}
  \proofstep{1}{\EqRel{R,A}}{\rA}
  \proofstep{1}{%
    \begin{aligned}[t]
      &\exists! q\colon A\to\QuotientSet{A}{R}\,{}\\
      &\forall x\in A\,q(x)=\tQuotProj{A}{R}(x)
    \end{aligned}}{%
    \FormulaRefAuto{%
      x\in A\vdash t(x)\coloneq y\dsep x\in A\vdash t(x)\in B
      \exists! F\colon A\to B\,\forall x\in A\,F(x)=t(x)
    }{%
      \FormulaRefAuto{QuotProjRuleDef}{1},
      \FormulaRefAuto{QuotProjRuleInQuotient}{1}}}
\end{tabproof}

\FormulaDefDeltaK[Quotientenprojektion]{%
  \EqRel{R,A}\vdash
  \QuotProj{A}{R}\coloneqq
  \iota q\Bigl(q\colon A\to\QuotientSet{A}{R}\land
  \forall x\in A\,q(x)=\tQuotProj{A}{R}(x)\Bigr)
}{QuotProjDef}{%
  \DeltaRow{Mengen}{A\dsep R\dsep x}
  \DeltaRow{Funktionen}{q\dsep \QuotProj{A}{R}}
  \DeltaRow{Funktionensymbol}{\tQuotProj{A}{R}}
  \DeltaRow{Neue Symbole}{\QuotProj{A}{R}}
}

\FormulaThmDeltaK[Typisierung der Quotientenprojektion]{%
  \EqRel{R,A}\vdash
  \QuotProj{A}{R}\colon A\to\QuotientSet{A}{R}
}{QuotProjFunction}{%
  \DeltaRow{Mengen}{A\dsep R}
  \DeltaRow{Funktionen}{\QuotProj{A}{R}}
}
\begin{tabproof}
  \proofstep{1}{\EqRel{R,A}}{\rA}
  \proofstep{1}{\QuotProj{A}{R}\colon A\to\QuotientSet{A}{R}}{%
    \rAEa{\FormulaRefAuto{QuotProjDef}{1}}}
\end{tabproof}

\FormulaThmDeltaK[Wert der Quotientenprojektion]{%
  \EqRel{R,A}\dsep x\in A\vdash
  \QuotProj{A}{R}(x)=\EqClass{x}{R}{A}
}{QuotProjValue}{%
  \DeltaRow{Mengen}{A\dsep R\dsep x}
  \DeltaRow{Funktionen}{\QuotProj{A}{R}}
  \DeltaRow{Äquivalenzklassen}{\EqClass{x}{R}{A}}
}
\begin{tabproof}
  \proofstep{1}{\EqRel{R,A}}{\rA}
  \proofstep{2}{x\in A}{\rA}
  \proofstep{1}{\forall x\in A\,
    \QuotProj{A}{R}(x)=\tQuotProj{A}{R}(x)}{%
    \rAEb{\FormulaRefAuto{QuotProjDef}{1}}}
  \proofstep{1,2}{\QuotProj{A}{R}(x)=\tQuotProj{A}{R}(x)}{%
    \rRE{2,\rUE{3}}}
  \proofstep{1,2}{\tQuotProj{A}{R}(x)=\EqClass{x}{R}{A}}{%
    \FormulaRefAuto{QuotProjRuleDef}{1,2}}
  \proofstep{1,2}{\QuotProj{A}{R}(x)=\EqClass{x}{R}{A}}{%
    \FormulaRefAuto{a = b,\, b = c \vdash a = c}{4,5}}
\end{tabproof}

\FormulaDefDeltaK[Quotientenbildabbildung einer Teilfamilie]{%
  \EqRel{R,A}\dsep M\subseteq\powerset(A)\vdash
  \QuotFamMap{A}{R}{M}\coloneqq
  \PowMap{\QuotProj{A}{R}}\restriction_M
}{QuotFamMapDef}{%
  \DeltaRow{Mengen}{A\dsep R\dsep M}
  \DeltaRow{Funktionen}{\QuotProj{A}{R}\dsep \PowMap{\QuotProj{A}{R}}\dsep
  \QuotFamMap{A}{R}{M}}
  \DeltaRow{Neue Symbole}{\QuotFamMap{A}{R}{M}}
}

\FormulaThmDeltaK[Typisierung der Quotientenbildabbildung]{%
  \EqRel{R,A}\dsep M\subseteq\powerset(A)\vdash
  \QuotFamMap{A}{R}{M}\colon
  M\to\powerset(\QuotientSet{A}{R})
}{QuotFamMapFunction}{%
  \DeltaRow{Mengen}{A\dsep R\dsep M}
  \DeltaRow{Funktionen}{\QuotProj{A}{R}\dsep \QuotFamMap{A}{R}{M}}
}
\begin{tabproof}
  \proofstep{1}{\EqRel{R,A}}{\rA}
  \proofstep{2}{M\subseteq\powerset(A)}{\rA}
  \proofstep{1}{\QuotProj{A}{R}\colon A\to\QuotientSet{A}{R}}{%
    \FormulaRefAuto{QuotProjFunction}{1}}
  \proofstep{1}{\PowMap{\QuotProj{A}{R}}\colon
    \powerset(A)\to\powerset(\QuotientSet{A}{R})}{%
    \FormulaRefAuto{F\colon A\to B \vdash
    \PowMap{F}\colon \powerset(A)\to \powerset(B)}{3}}
  \proofstep{1,2}{\PowMap{\QuotProj{A}{R}}\restriction_M
    \colon M\to\powerset(\QuotientSet{A}{R})}{%
    \FormulaRefAuto{C\subseteq A \vdash F\restriction_C\colon C\to B}{2,4}}
  \proofstep{1,2}{\QuotFamMap{A}{R}{M}=
    \PowMap{\QuotProj{A}{R}}\restriction_M}{%
    \FormulaRefAuto{QuotFamMapDef}{1,2}}
  \proofstep{1,2}{\QuotFamMap{A}{R}{M}\colon
    M\to\powerset(\QuotientSet{A}{R})}{\rIE{6,5}}
\end{tabproof}

\FormulaThmDeltaK[Wert der Quotientenbildabbildung]{%
  \EqRel{R,A}\dsep M\subseteq\powerset(A)\dsep X\in M\vdash
  \QuotFamMap{A}{R}{M}(X)=\QuotProj{A}{R}[X]
}{QuotFamMapValue}{%
  \DeltaRow{Mengen}{A\dsep R\dsep M\dsep X}
  \DeltaRow{Funktionen}{\QuotProj{A}{R}\dsep \QuotFamMap{A}{R}{M}}
}
\begin{tabproof}
  \proofstep{1}{\EqRel{R,A}}{\rA}
  \proofstep{2}{M\subseteq\powerset(A)}{\rA}
  \proofstep{3}{X\in M}{\rA}
  \proofstep{1}{\QuotProj{A}{R}\colon A\to\QuotientSet{A}{R}}{%
    \FormulaRefAuto{QuotProjFunction}{1}}
  \proofstep{1}{\PowMap{\QuotProj{A}{R}}\colon
    \powerset(A)\to\powerset(\QuotientSet{A}{R})}{%
    \FormulaRefAuto{F\colon A\to B \vdash
    \PowMap{F}\colon \powerset(A)\to \powerset(B)}{4}}
  \proofstep{1,2}{\QuotFamMap{A}{R}{M}=
    \PowMap{\QuotProj{A}{R}}\restriction_M}{%
    \FormulaRefAuto{QuotFamMapDef}{1,2}}
  \proofstep{1,2,3}{\QuotFamMap{A}{R}{M}(X)=
    \bigl(\PowMap{\QuotProj{A}{R}}\restriction_M\bigr)(X)}{%
    \FormulaRefAuto{x\in A\dsep F=G\vdash F(x)=G(x)}{3,6}}
  \proofstep{1,2,3}{\bigl(\PowMap{\QuotProj{A}{R}}\restriction_M\bigr)(X)=
    \PowMap{\QuotProj{A}{R}}(X)}{%
    \FormulaRefAuto{C\subseteq A \dsep x\in C \vdash F\restriction_C(x)=F(x)}{2,3,5}}
  \proofstep{2,3}{X\in\powerset(A)}{%
    \FormulaRefAuto{A\subseteq B,\, x\in A \vdash x\in B}{2,3}}
  \proofstep{1,2,3}{\PowMap{\QuotProj{A}{R}}(X)=
    \QuotProj{A}{R}[X]}{%
    \FormulaRefAuto{F\colon A\to B \dsep X\in\powerset(A)\vdash
    \PowMap{F}(X)=F[X]}{4,9}}
  \proofstep{1,2,3}{\QuotFamMap{A}{R}{M}(X)=
    \QuotProj{A}{R}[X]}{%
    \FormulaRefAuto{a = b,\, b = c \vdash a = c}{7,
    \FormulaRefAuto{a = b,\, b = c \vdash a = c}{8,10}}}
\end{tabproof}

\subsubsection{Charakterisierung der Quotientenbildabbildung}

\FormulaThmDeltaK[Quotientenprojektionsbilder liegen in der Quotientenmenge]{%
  \EqRel{R,A}\dsep M\subseteq\powerset(A)\dsep X\in M
  \dsep C\in\QuotProj{A}{R}[X]
  \vdash
  C\in\QuotientSet{A}{R}
}{QuotProjImageInQuotient}{%
  \DeltaRow{Mengen}{A\dsep R\dsep M\dsep X\dsep C\dsep a}
  \DeltaRow{Funktionen}{\QuotProj{A}{R}}
  \DeltaRow{Quotientenmenge}{\QuotientSet{A}{R}}
}
\begin{tabproof}
  \proofstep{1}{\EqRel{R,A}}{\rA}
  \proofstep{2}{M\subseteq\powerset(A)}{\rA}
  \proofstep{3}{X\in M}{\rA}
  \proofstep{4}{C\in\QuotProj{A}{R}[X]}{\rA}
  \proofstep{2,3}{X\in\powerset(A)}{%
    \FormulaRefAuto{A\subseteq B,\, x\in A \vdash x\in B}{2,3}}
  \proofstep{2,3}{X\subseteq A}{%
    \FormulaRefAuto{x \in \powerset(A)\;\eqvdash\; x \subseteq A}{5}}
  \proofstep{1}{\QuotProj{A}{R}\colon
    A\to\QuotientSet{A}{R}}{%
    \FormulaRefAuto{QuotProjFunction}{1}}
  \proofstep{1,2,3}{\QuotProj{A}{R}[X]\subseteq
    \QuotientSet{A}{R}}{%
    \FormulaRefAuto{C\subseteq A \dsep F\colon A\to B \vdash F[C]\subseteq B}{6,7}}
  \proofstep{1,2,3,4}{C\in\QuotientSet{A}{R}}{%
    \FormulaRefAuto{A\subseteq B,\, x\in A \vdash x\in B}{8,4}}
\end{tabproof}

\FormulaThmDeltaK[Quotientenprojektionsbilder liefern Klassenzeugen]{%
  \EqRel{R,A}\dsep M\subseteq\powerset(A)\dsep X\in M
  \dsep C\in\QuotProj{A}{R}[X]
  \vdash
  \exists a\in X\,C=\EqClass{a}{R}{A}
}{QuotProjImageClassWitness}{%
  \DeltaRow{Mengen}{A\dsep R\dsep M\dsep X\dsep C\dsep a}
  \DeltaRow{Funktionen}{\QuotProj{A}{R}}
  \DeltaRow{Quotientenmenge}{\QuotientSet{A}{R}}
}
\begin{tabproofwide}
  \proofstepwidestar[1]{\EqRel{R,A}}{\rA}
  \proofstepwidestar[2]{M\subseteq\powerset(A)}{\rA}
  \proofstepwidestar[3]{X\in M}{\rA}
  \proofstepwidestar[4]{C\in\QuotProj{A}{R}[X]}{\rA}
  \proofstepwidestar[2,3]{X\in\powerset(A)}{%
    \FormulaRefAuto{A\subseteq B,\, x\in A \vdash x\in B}{2,3}}
  \proofstepwidestar[2,3]{X\subseteq A}{%
    \FormulaRefAuto{x \in \powerset(A)\;\eqvdash\; x \subseteq A}{5}}
  \proofstepwidestar[1]{\QuotProj{A}{R}\colon
    A\to\QuotientSet{A}{R}}{%
    \FormulaRefAuto{QuotProjFunction}{1}}
  \proofstepwidestar[1,2,3,4]{\exists a\in X\,
    C=\QuotProj{A}{R}(a)}{%
    \FormulaRefAuto{C\subseteq A\dsep F\colon A\to B\dsep
    y\in F[C]\vdash \exists x\in C\,y=F(x)}{6,7,4}}
  \proofstepwidestar[9]{a\in X\land C=\QuotProj{A}{R}(a)}{\rA}
  \proofstepwidestar[9]{a\in X}{\rAEa{9}}
  \proofstepwidestar[9]{C=\QuotProj{A}{R}(a)}{\rAEb{9}}
  \proofstepwidestar[2,3,9]{a\in A}{%
    \FormulaRefAuto{A\subseteq B,\, x\in A \vdash x\in B}{6,10}}
  \proofstepwidestar[1,9]{\QuotProj{A}{R}(a)=\EqClass{a}{R}{A}}{%
    \FormulaRefAuto{QuotProjValue}{1,12}}
  \proofstepwidestar[1,9]{C=\EqClass{a}{R}{A}}{%
    \FormulaRefAuto{a = b,\, b = c \vdash a = c}{11,13}}
  \proofstepwidestar[1,9]{a\in X\land C=\EqClass{a}{R}{A}}{\rAI{10,14}}
  \proofstepwidestar[1,9]{\exists a\in X\,C=\EqClass{a}{R}{A}}{\rEI{15}}
  \proofstepwidestar[1,2,3,4]{\exists a\in X\,C=\EqClass{a}{R}{A}}{%
    \rEE{8,9,16}}
\end{tabproofwide}

\FormulaThmDeltaK[Quotientenprojektionsbilder liefern Klassen]{%
  \EqRel{R,A}\dsep M\subseteq\powerset(A)\dsep X\in M
  \dsep C\in\QuotProj{A}{R}[X]
  \vdash
  C\in\QuotientSet{A}{R}\land
  \exists a\in X\,C=\EqClass{a}{R}{A}
}{QuotProjImageToEqClass}{%
  \DeltaRow{Mengen}{A\dsep R\dsep M\dsep X\dsep C\dsep a}
  \DeltaRow{Funktionen}{\QuotProj{A}{R}}
  \DeltaRow{Quotientenmenge}{\QuotientSet{A}{R}}
}
\begin{tabproof}
  \proofstep{1}{\EqRel{R,A}}{\rA}
  \proofstep{2}{M\subseteq\powerset(A)}{\rA}
  \proofstep{3}{X\in M}{\rA}
  \proofstep{4}{C\in\QuotProj{A}{R}[X]}{\rA}
  \proofstep{1,2,3,4}{C\in\QuotientSet{A}{R}}{%
    \FormulaRefAuto{QuotProjImageInQuotient}{1,2,3,4}}
  \proofstep{1,2,3,4}{\exists a\in X\,C=\EqClass{a}{R}{A}}{%
    \FormulaRefAuto{QuotProjImageClassWitness}{1,2,3,4}}
  \proofstep{1,2,3,4}{C\in\QuotientSet{A}{R}\land
    \exists a\in X\,C=\EqClass{a}{R}{A}}{\rAI{5,6}}
\end{tabproof}

\FormulaThmDeltaK[Klassen liefern Quotientenprojektionsbilder]{%
  \EqRel{R,A}\dsep M\subseteq\powerset(A)\dsep X\in M
  \dsep \exists a\in X\,C=\EqClass{a}{R}{A}
  \vdash
  C\in\QuotProj{A}{R}[X]
}{EqClassToQuotProjImage}{%
  \DeltaRow{Mengen}{A\dsep R\dsep M\dsep X\dsep C\dsep a}
  \DeltaRow{Funktionen}{\QuotProj{A}{R}}
  \DeltaRow{Quotientenmenge}{\QuotientSet{A}{R}}
}
\begin{tabproofwide}
  \proofstepwidestar[1]{\EqRel{R,A}}{\rA}
  \proofstepwidestar[2]{M\subseteq\powerset(A)}{\rA}
  \proofstepwidestar[3]{X\in M}{\rA}
  \proofstepwidestar[4]{\exists a\in X\,C=\EqClass{a}{R}{A}}{\rA}
  \proofstepwidestar[2,3]{X\in\powerset(A)}{%
    \FormulaRefAuto{A\subseteq B,\, x\in A \vdash x\in B}{2,3}}
  \proofstepwidestar[2,3]{X\subseteq A}{%
    \FormulaRefAuto{x \in \powerset(A)\;\eqvdash\; x \subseteq A}{5}}
  \proofstepwidestar[1]{\QuotProj{A}{R}\colon
    A\to\QuotientSet{A}{R}}{%
    \FormulaRefAuto{QuotProjFunction}{1}}
  \proofstepwidestar[8]{a\in X\land C=\EqClass{a}{R}{A}}{\rA}
  \proofstepwidestar[8]{a\in X}{\rAEa{8}}
  \proofstepwidestar[8]{C=\EqClass{a}{R}{A}}{\rAEb{8}}
  \proofstepwidestar[2,3,8]{a\in A}{%
    \FormulaRefAuto{A\subseteq B,\, x\in A \vdash x\in B}{6,9}}
  \proofstepwidestar[1,8]{\QuotProj{A}{R}(a)=\EqClass{a}{R}{A}}{%
    \FormulaRefAuto{QuotProjValue}{1,11}}
  \proofstepwidestar[1,8]{\EqClass{a}{R}{A}=\QuotProj{A}{R}(a)}{%
    \FormulaRefAuto{a=b\vdash b=a}{12}}
  \proofstepwidestar[1,8]{C=\QuotProj{A}{R}(a)}{%
    \FormulaRefAuto{a = b,\, b = c \vdash a = c}{10,13}}
  \proofstepwidestar[1,8]{a\in X\land C=\QuotProj{A}{R}(a)}{%
    \rAI{9,14}}
  \proofstepwidestar[1,8]{\exists a\in X\,
    C=\QuotProj{A}{R}(a)}{\rEI{15}}
  \proofstepwidestar[1,2,3,8]{C\in\QuotProj{A}{R}[X]}{%
    \FormulaRefAuto{C\subseteq A\dsep F\colon A\to B\dsep
    \exists x\in C\,y=F(x)\vdash y\in F[C]}{6,7,16}}
  \proofstepwidestar[1,2,3,4]{C\in\QuotProj{A}{R}[X]}{\rEE{4,8,17}}
\end{tabproofwide}

\FormulaThmDeltaK[Charakterisierung der Quotientenbildabbildung]{%
  \EqRel{R,A}\dsep M\subseteq\powerset(A)\dsep X\in M\vdash
  \begin{aligned}[t]
    C\in\QuotFamMap{A}{R}{M}(X)
    \leftrightarrow{}&
    C\in\QuotientSet{A}{R}\land{}\\
    &\exists a\in X\,C=\EqClass{a}{R}{A}
  \end{aligned}
}{QuotFamMapChar}{%
  \DeltaRow{Mengen}{A\dsep R\dsep M\dsep X\dsep C\dsep a}
  \DeltaRow{Funktionen}{\QuotProj{A}{R}\dsep \QuotFamMap{A}{R}{M}}
  \DeltaRow{Quotientenmenge}{\QuotientSet{A}{R}}
}
\begin{tabproofsplitwide}
  \proofpartwideindR[\(\rightarrow\)]{%
    \EqRel{R,A}\dsep M\subseteq\powerset(A)\dsep X\in M
    \vdash
    C\in\QuotFamMap{A}{R}{M}(X)\rightarrow
    \bigl(C\in\QuotientSet{A}{R}\land
    \exists a\in X\,C=\EqClass{a}{R}{A}\bigr)
  }{%
    \EqRel{R,A}\dsep M\subseteq\powerset(A)\dsep X\in M
    \vdash
    C\in\QuotFamMap{A}{R}{M}(X)\rightarrow
    \bigl(C\in\QuotientSet{A}{R}\land
    \exists a\in X\,C=\EqClass{a}{R}{A}\bigr)
  }[QuotFamMapCharForward]
    \proofstepwidestar[1]{\EqRel{R,A}}{\rA}
    \proofstepwidestar[2]{M\subseteq\powerset(A)}{\rA}
    \proofstepwidestar[3]{X\in M}{\rA}
    \proofstepwidestar[4]{C\in\QuotFamMap{A}{R}{M}(X)}{\rA}
    \proofstepwidestar[1,2,3]{\QuotFamMap{A}{R}{M}(X)=
      \QuotProj{A}{R}[X]}{%
      \FormulaRefAuto{QuotFamMapValue}{1,2,3}}
    \proofstepwidestar[1,2,3,4]{C\in\QuotProj{A}{R}[X]}{\rIE{5,4}}
    \proofstepwidestar[1,2,3,4]{C\in\QuotientSet{A}{R}\land
      \exists a\in X\,C=\EqClass{a}{R}{A}}{%
      \FormulaRefAuto{QuotProjImageToEqClass}{1,2,3,6}}
    \proofstepwidestar[1,2,3]{C\in\QuotFamMap{A}{R}{M}(X)
      \rightarrow
      \bigl(C\in\QuotientSet{A}{R}\land
      \exists a\in X\,C=\EqClass{a}{R}{A}\bigr)}{\rRI{4,7}}
  \closeproofpartwideindR

  \proofpartwideindR[\(\leftarrow\)]{%
    \EqRel{R,A}\dsep M\subseteq\powerset(A)\dsep X\in M
    \vdash
    \bigl(C\in\QuotientSet{A}{R}\land
    \exists a\in X\,C=\EqClass{a}{R}{A}\bigr)
    \rightarrow C\in\QuotFamMap{A}{R}{M}(X)
  }{%
    \EqRel{R,A}\dsep M\subseteq\powerset(A)\dsep X\in M
    \vdash
    \bigl(C\in\QuotientSet{A}{R}\land
    \exists a\in X\,C=\EqClass{a}{R}{A}\bigr)
    \rightarrow C\in\QuotFamMap{A}{R}{M}(X)
  }[QuotFamMapCharBackward]
    \proofstepwidestar[1]{\EqRel{R,A}}{\rA}
    \proofstepwidestar[2]{M\subseteq\powerset(A)}{\rA}
    \proofstepwidestar[3]{X\in M}{\rA}
    \proofstepwidestar[4]{C\in\QuotientSet{A}{R}\land
      \exists a\in X\,C=\EqClass{a}{R}{A}}{\rA}
    \proofstepwidestar[4]{\exists a\in X\,C=\EqClass{a}{R}{A}}{\rAEb{4}}
    \proofstepwidestar[1,2,3,4]{C\in\QuotProj{A}{R}[X]}{%
      \FormulaRefAuto{EqClassToQuotProjImage}{1,2,3,5}}
    \proofstepwidestar[1,2,3]{\QuotProj{A}{R}[X]=
      \QuotFamMap{A}{R}{M}(X)}{%
      \FormulaRefAuto{a=b\vdash b=a}{%
      \FormulaRefAuto{QuotFamMapValue}{1,2,3}}}
    \proofstepwidestar[1,2,3,4]{C\in\QuotFamMap{A}{R}{M}(X)}{%
      \rIE{7,6}}
    \proofstepwidestar[1,2,3]{%
      \bigl(C\in\QuotientSet{A}{R}\land
      \exists a\in X\,C=\EqClass{a}{R}{A}\bigr)
      \rightarrow C\in\QuotFamMap{A}{R}{M}(X)}{\rRI{4,8}}
  \closeproofpartwideindR

  \proofpartwide{Schluss}
    \proofstepwidestar[1]{\EqRel{R,A}}{\rA}
    \proofstepwidestar[2]{M\subseteq\powerset(A)}{\rA}
    \proofstepwidestar[3]{X\in M}{\rA}
    \proofstepwidestar[1,2,3]{C\in\QuotFamMap{A}{R}{M}(X)
      \rightarrow
      \bigl(C\in\QuotientSet{A}{R}\land
      \exists a\in X\,C=\EqClass{a}{R}{A}\bigr)}{%
      \FormulaRefAuto{QuotFamMapCharForward}{1,2,3}}
    \proofstepwidestar[1,2,3]{%
      \bigl(C\in\QuotientSet{A}{R}\land
      \exists a\in X\,C=\EqClass{a}{R}{A}\bigr)
      \rightarrow C\in\QuotFamMap{A}{R}{M}(X)}{%
      \FormulaRefAuto{QuotFamMapCharBackward}{1,2,3}}
    \proofstepwidestar[1,2,3]{C\in\QuotFamMap{A}{R}{M}(X)
      \leftrightarrow
      \bigl(C\in\QuotientSet{A}{R}\land
      \exists a\in X\,C=\EqClass{a}{R}{A}\bigr)}{\rLRI{4,5}}
  \closeproofpartwide
\end{tabproofsplitwide}


\section{Gleichmächtigkeit}

\FormulaDefDeltaK[Begriff der gleichmächtigen Mengen]{A \EqCard B \coloneqq \exists F\colon A\bij B}{Gleichmächtigkeit}{
  \DeltaRow{Mengen}{A\dsep B}
  \DeltaRow{\textbf{Neue Symbole}}{}
  \DeltaRow{Gleichmächtigkeit}{}
}

\FormulaThmDeltaK[Frische Adjunktion erhält Gleichmächtigkeit]{%
  A\EqCard B\dsep a\notin A\dsep b\notin B
  \vdash A\cup\{a\}\EqCard B\cup\{b\}%
}{EqCardFreshAdjunction}{%
  \DeltaRow{Mengen}{A\dsep B\dsep a\dsep b}%
  \DeltaRow{Funktionen}{F}%
}
\begin{tabproof}
  \proofstep{1}{A\EqCard B}{\rA}
  \proofstep{2}{a\notin A}{\rA}
  \proofstep{3}{b\notin B}{\rA}
  \proofstep{1}{\exists F\colon A\bij B}{%
    \FormulaRefAuto{A \EqCard B \coloneqq \exists F\colon A\bij B}{1}}
  \proofstep{5}{F\colon A\bij B}{\rA}
  \proofstep{2,3,5}{%
    \ExtMap{\iota_{B,B\cup\{b\}}\circ F}{a}{b}
    \colon A\cup\{a\}\bij B\cup\{b\}}{%
    \FormulaRefAuto{FreshExtMapBijective}[thm]{5,2,3}}
  \proofstep{2,3,5}{\exists G\colon A\cup\{a\}\bij B\cup\{b\}}{\rEI{6}}
  \proofstep{2,3,5}{A\cup\{a\}\EqCard B\cup\{b\}}{%
    \FormulaRefAuto{A \EqCard B \coloneqq \exists F\colon A\bij B}{7}}
  \proofstep{1,2,3}{A\cup\{a\}\EqCard B\cup\{b\}}{\rEE{4,5,8}}
\end{tabproof}

\FormulaThmDeltaK[Reflexivität der Gleichmächtigkeit]{%
  A \EqCard A%
}{EqCardReflexive}{
  \DeltaRow{Mengen}{A}
}
\begin{tabproof}
  \proofstep{}{\Id_A\colon A\bij A}{\FormulaRefAuto{\Id_A\colon A\bij A}}
  \proofstep{}{\exists F\colon A\bij A}{\rEI{1}}
  \proofstep{}{A \EqCard A}{\FormulaRefAuto{A \EqCard B \coloneqq \exists F\colon A\bij B}{2}}
\end{tabproof}

\FormulaThmDeltaK[Symmetrie der Gleichmächtigkeit]{%
  A \EqCard B \vdash B \EqCard A%
}{EqCardSymmetric}{
  \DeltaRow{Mengen}{A\dsep B}
  \DeltaRow{Funktionen}{F\dsep G}
}
\begin{tabproof}
  \proofstep{1}{A \EqCard B}{\rA}
  \proofstep{1}{\exists F\colon A\bij B}{\FormulaRefAuto{A \EqCard B \coloneqq \exists F\colon A\bij B}{1}}

  \proofstep{3}{F\colon A\bij B}{\rA}
  \proofstep{3}{F^{-1}\colon B\bij A}{\FormulaRefAuto{F\colon A\bij B\vdash F^{-1}\colon B\bij A}{3}}
  \proofstep{3}{\exists G\colon B\bij A}{\rEI{4}}

  \proofstep{1}{\exists G\colon B\bij A}{\rEE{2,3,5}}
  \proofstep{1}{B \EqCard A}{\FormulaRefAuto{A \EqCard B \coloneqq \exists F\colon A\bij B}{6}}
\end{tabproof}

\FormulaThmDeltaK[Transitivität der Gleichmächtigkeit]{%
  A \EqCard B \dsep B \EqCard C \vdash A \EqCard C%
}{EqCardTransitive}{
  \DeltaRow{Mengen}{A\dsep B\dsep C}
  \DeltaRow{Funktionen}{F\dsep G\dsep H}
}
\begin{tabproof}
  \proofstep{1}{A \EqCard B}{\rA}
  \proofstep{2}{B \EqCard C}{\rA}

  \proofstep{1}{\exists F\colon A\bij B}{\FormulaRefAuto{A \EqCard B \coloneqq \exists F\colon A\bij B}{1}}
  \proofstep{2}{\exists G\colon B\bij C}{\FormulaRefAuto{A \EqCard B \coloneqq \exists F\colon A\bij B}{2}}

  \proofstep{5}{F\colon A\bij B}{\rA}
  \proofstep{6}{G\colon B\bij C}{\rA}
  \proofstep{5,6}{G\circ F\colon A\bij C}{\FormulaRefAuto{F\colon A\bij B \dsep G\colon B\bij C\vdash G\circ F\colon A\bij C}{5,6}}
  \proofstep{5,6}{\exists H\colon A\bij C}{\rEI{7}}
  \proofstep{5}{\exists H\colon A\bij C}{\rEE{4,6,8}}
  \proofstep{1,2}{\exists H\colon A\bij C}{\rEE{3,5,9}}

  \proofstep{1,2}{A \EqCard C}{\FormulaRefAuto{A \EqCard B \coloneqq \exists F\colon A\bij B}{10}}
\end{tabproof}

\FormulaThmDeltaK[Gleichmächtigkeit ist eine Äquivalenzrelation]{%
  \begin{aligned}[t]
    \vdash{}& A\EqCard A\land{}\\
    &\bigl(A\EqCard B\rightarrow B\EqCard A\bigr)\land{}\\
    &\bigl((A\EqCard B\land B\EqCard C)\rightarrow A\EqCard C\bigr)
  \end{aligned}
}{EqCardEquivalence}{%
  \DeltaRow{Mengen}{A\dsep B\dsep C}
}
\begin{tabproofwide}
  \proofstepwidestar[]{A\EqCard A}{%
    \FormulaRefAuto{EqCardReflexive}}
  \proofstepwidestar[2]{A\EqCard B}{\rA}
  \proofstepwidestar[2]{B\EqCard A}{%
    \FormulaRefAuto{EqCardSymmetric}{2}}
  \proofstepwidestar[]{A\EqCard B\rightarrow B\EqCard A}{%
    \rRI{2,3}}
  \proofstepwidestar[5]{A\EqCard B\land B\EqCard C}{\rA}
  \proofstepwidestar[5]{A\EqCard B}{\rAEa{5}}
  \proofstepwidestar[5]{B\EqCard C}{\rAEb{5}}
  \proofstepwidestar[5]{A\EqCard C}{%
    \FormulaRefAuto{EqCardTransitive}{6,7}}
  \proofstepwidestar[]{%
    (A\EqCard B\land B\EqCard C)\rightarrow A\EqCard C}{%
    \rRI{5,8}}
  \proofstepwidestar[]{%
    A\EqCard A\land
    \bigl(A\EqCard B\rightarrow B\EqCard A\bigr)\land
    \bigl((A\EqCard B\land B\EqCard C)\rightarrow A\EqCard C\bigr)}{%
    \rAI{1,\rAI{4,9}}}
\end{tabproofwide}

\FormulaThmDeltaK[Transport von Injektionen entlang der Gleichmächtigkeit]{%
  A \EqCard M\dsep B \EqCard N\dsep F\colon A\inj B
  \vdash \exists G\colon M\inj N
}{EqCardInjectionTransport}{%
  \DeltaRow{Mengen}{A\dsep B\dsep M\dsep N}
  \DeltaRow{Funktionen}{F\dsep G\dsep H\dsep K}
}
\begin{tabproof}
  \proofstep{1}{A \EqCard M}{\rA}
  \proofstep{2}{B \EqCard N}{\rA}
  \proofstep{3}{F\colon A\inj B}{\rA}
  \proofstep{1}{\exists H\colon A\bij M}{%
    \FormulaRefAuto{A \EqCard B \coloneqq \exists F\colon A\bij B}{1}}
  \proofstep{2}{\exists K\colon B\bij N}{%
    \FormulaRefAuto{A \EqCard B \coloneqq \exists F\colon A\bij B}{2}}

  \proofstep{6}{H\colon A\bij M}{\rA}
  \proofstep{7}{K\colon B\bij N}{\rA}
  \proofstep{6}{H^{-1}\colon M\bij A}{%
    \FormulaRefAuto{F\colon A\bij B\vdash F^{-1}\colon B\bij A}{6}}
  \proofstep{6}{H^{-1}\colon M\inj A}{%
    \FormulaRefAuto{F\colon A\bij B \vdash F\colon A\inj B}{8}}
  \proofstep{3,6}{F\circ H^{-1}\colon M\inj B}{%
    \FormulaRefAuto{F\colon A\inj B \dsep G\colon B\inj C\vdash G\circ F\colon A\inj C}{9,3}}
  \proofstep{7}{K\colon B\inj N}{%
    \FormulaRefAuto{F\colon A\bij B \vdash F\colon A\inj B}{7}}
  \proofstep{3,6,7}{K\circ(F\circ H^{-1})\colon M\inj N}{%
    \FormulaRefAuto{F\colon A\inj B \dsep G\colon B\inj C\vdash G\circ F\colon A\inj C}{10,11}}
  \proofstep{3,6,7}{\exists G\colon M\inj N}{\rEI{12}}

  \proofstep{2,3,6}{\exists G\colon M\inj N}{\rEE{5,7,13}}
  \proofstep{1,2,3}{\exists G\colon M\inj N}{\rEE{4,6,14}}
\end{tabproof}


\FormulaThmDelta[Potenzmengen gleichmächtiger Mengen sind gleichmächtig]{%
  A \EqCard B \vdash \powerset(A)\EqCard \powerset(B)
}{%
  \DeltaRow{Mengen}{A\dsep B}
  \DeltaRow{Funktionen}{F\dsep G\dsep \widehat{F}}
}
\begin{tabproof}
  \proofstep{1}{A \EqCard B}{\rA}
  \proofstep{1}{\exists F\colon A\bij B}{%
    \FormulaRefAuto{A \EqCard B \coloneqq \exists F\colon A\bij B}{1}}

  \proofstep{3}{F\colon A\bij B}{\rA}
    \proofstep{3}{\PowMap{F}\colon \powerset(A)\bij \powerset(B)}{%
      \FormulaRefAuto{F\colon A\bij B \vdash \PowMap{F}\colon \powerset(A)\bij \powerset(B)}{3}}
  \proofstep{3}{\exists G\colon \powerset(A)\bij \powerset(B)}{\rEI{4}}

  \proofstep{1}{\exists G\colon \powerset(A)\bij \powerset(B)}{\rEE{2,3,5}}
  \proofstep{1}{\powerset(A)\EqCard \powerset(B)}{%
    \FormulaRefAuto{A \EqCard B \coloneqq \exists F\colon A\bij B}{6}}
\end{tabproof}


\end{document}
